% ============================================================================
% 附录：评分标准全文、劣化算子、数据结构、闸门、角色提示、第三方组件、复现说明
% 依据：eval/rubrics/report_v2.json、finagent/evaluation/calibrate.py、data/extractions/、finagent/output/sources.py、
%       finagent/report_pipeline.py、.claude/skills/、THIRD_PARTY.md、README.md、CLAUDE.md、scripts/
% ============================================================================

\section{评分标准全文}
\label{app:rubric}

本附录列出默认评分标准 \file{eval/rubrics/report\_v2.json} 的全部 41 个条目（正文见第~\ref{sec:eval-rubric}~节）。
“判定”列中，程序条目附检查函数名（\file{finagent/evaluation/checks.py} 中的 \code{check\_<名>}）；
“路由”列为未通过时问题单回传的会话（\code{generator} 为生成侧，\code{reviewer} 为复核侧；阶段以 10 开头的复核侧问题交给排版设计师）与对应阶段。
交给盲评会话的 \file{rubric.json} 只含 id、条目、判定标准、反例与是否负向，不含本表的分值、层级与路由。

\paragraph{全局参数。}
\code{gate\_cap} $=40$（L0 任一失败的总分上限）；\code{negative\_floor} $=-20$（负向扣分合计下限）；
交付标准 \code{delivery}：总分 $\ge 75$，L1 / L2 / L3 得分率分别 $\ge 80\%$ / $60\%$ / $60\%$，\code{allow\_negative} $=$ false（不得命中负向条目）。
各层分值合计：L0 为 0，L1 为 35，L2 为 45，L3 为 20，N 为 $-30$。

{\footnotesize
\setlength{\tabcolsep}{3pt}
\begin{longtable}{@{}p{0.07\linewidth}p{0.13\linewidth}p{0.12\linewidth}p{0.07\linewidth}p{0.12\linewidth}p{0.43\linewidth}@{}}
\caption{评分标准 \code{report\_v2} 全部条目。}\label{tab:app-rubric}\\
\toprule
id & 条目 & 判定 & 分值 & 路由 / 阶段 & 判定标准（及反例） \\
\midrule
\endfirsthead
\toprule
id & 条目 & 判定 & 分值 & 路由 / 阶段 & 判定标准（及反例） \\
\midrule
\endhead
\bottomrule
\endfoot
L0-01 & 数字护栏 & 程序\newline\code{guard} & 0 & generator\newline 8 文稿 & 当前文稿的每个数字都能在本次运行产物或已登记外部来源中找到，指标与期间绑定无误。 \\
L0-02 & 三表勾稽 & 程序\newline\code{ts\_integrity} & 0 & generator\newline 7 估值运行 & 三表模型资产负债平衡、现金衔接、各项滚动全部通过，无硬性失败。 \\
L0-03 & DCF 与三表同口径 & 程序\newline\code{dcf\_tie} & 0 & generator\newline 7 估值运行 & DCF 的收入、EBIT、折旧摊销、资本开支、营运资本逐期取自三表模型。 \\
L0-04 & 免责声明 & 程序\newline\code{disclaimer} & 0 & reviewer\newline 10 排版 & 报告含“不构成投资建议”的免责声明。 \\
L0-05 & 评级与估值一致 & 程序\newline\code{rating} & 0 & generator\newline 7 估值运行 & 评级按规则由上行空间得出，正文写的评级与之相同。 \\
\midrule
L1-01 & 抽取校验 & 程序\newline\code{extraction} & 3 & generator\newline 2 财务抽取 & 抽取事实逐字可定位，无拒收、冲突与勾稽失败。 \\
L1-02 & 分部预测基期闭合 & 程序\newline\code{segment\_base} & 3 & generator\newline 6 估值假设 & 分部驱动预测的基期分部收入加其他收入等于报表营业收入。 \\
L1-03 & 股权价值桥加总 & 程序\newline\code{bridge} & 4 & generator\newline 7 估值运行 & 各情景企业价值 + 加回项 $-$ 扣减项 = 股权价值。 \\
L1-04 & 每股价值复算 & 程序\newline\code{per\_share} & 3 & generator\newline 7 估值运行 & 各情景股权价值 $\div$ 总股本 = 每股价值。 \\
L1-05 & 情景排序 & 程序\newline\code{scenario\_order} & 2 & generator\newline 6 估值假设 & 悲观 $\le$ 基准 $\le$ 乐观。 \\
L1-06 & 敏感性方向 & 程序\newline\code{sensitivity} & 2 & generator\newline 7 估值运行 & WACC 上升价值下降、永续增长上升价值上升。 \\
L1-07 & 终值占比 & 程序\newline\code{terminal\_share} & 2 & generator\newline 6 估值假设 & 基准情景终值占企业价值不超过阈值。（参数：\code{max}=0.8） \\
L1-08 & 预测与历史衔接 & 程序\newline\code{forecast\_}\newline\code{continuity} & 2 & generator\newline 6 估值假设 & 首个预测年毛利率与最近实际年相差不超过阈值（百分点），收入增速变化不超过阈值。（参数：\code{gross\_margin\_pp}=5，\code{growth\_pp}=30） \\
L1-09 & 正文估值与产物一致 & 程序\newline\code{text\_valuation} & 3 & generator\newline 8 文稿 & 正文写出基准每股价值（按展示精度与产物一致）与评级。 \\
L1-10 & 外部信息可追溯 & 程序\newline\code{sources} & 4 & generator\newline 1 资料搜集 & sources.json 结构校验通过；正文至少引用一处 \texttt{[S编号]}，引用的编号都已登记；只能从外部资料找到出处的数字，所在句子带对应来源的编号。 \\
L1-11 & 股价有日期与来源 & 程序\newline\code{price} & 2 & generator\newline 6 估值假设 & 参考股价写明日期与来源。 \\
L1-12 & 复核改写过数字锁 & 程序\newline\code{numlock} & 3 & reviewer\newline 9 复核改写 & 复核改写前后，改写稿没有新增或改动数字。 \\
L1-13 & 人工复核清单 & 程序\newline\code{review\_list} & 2 & reviewer\newline 9 复核改写 & review.md 至少有一条填写完整的复核事项（事项、为什么重要、当前处理、请核对），或明确写“无”；只有表头不算。 \\
\midrule
L2-01 & 核心结论 & 评审 & 4 & generator\newline 8 文稿 & 首页或摘要用一两句话给出明确的核心观点：评级、估值结论，以及两三条支撑理由。读完摘要能复述“为什么值这个价”。\newline\textit{反例}：摘要只罗列财务数据，没有判断；支撑理由是“行业龙头”“前景广阔”这类无法检验的话 \\
L2-02 & 预期差 & 评审 & 4 & generator\newline 8 文稿 & 说明本报告与市场或一致预期的分歧在哪里，或者说明当前股价隐含了什么假设、我们为何不同意。\newline\textit{反例}：通篇没有提市场怎么看；只说“市场低估了公司价值”，不说低估了哪个变量 \\
L2-03 & 商业模式 & 评审 & 3 & generator\newline 8 文稿 & 讲清公司卖什么、卖给谁、靠什么赚钱，以及竞争优势体现在哪个可观察的指标上（份额、毛利率、成本、客户结构等）。\newline\textit{反例}：只有公司简介与发展历程；优势只用形容词描述 \\
L2-04 & 驱动拆解 & 评审 & 5 & generator\newline 8 文稿 & 至少一个核心驱动给出两环以上的定量推导：外部驱动 $\to$ 公司量、价或结构 $\to$ 收入或利润，每一环都有数字。\newline\textit{反例}：“受益于需求增长，收入有望提升”；只给出收入增速，不拆量与价 \\
L2-05 & 历史业绩归因 & 评审 & 4 & generator\newline 8 文稿 & 最近一年利润变动被归因到具体科目或驱动（毛利、费用、投资收益、减值等），并给出各自的量。\newline\textit{反例}：只写“净利润同比增长 X\%”；归因只有定性原因 \\
L2-06 & 行业与竞争 & 评审 & 4 & generator\newline 8 文稿 & 行业供需、周期位置或竞争格局有具体数据支撑（规模、增速、份额、价格等），且数据带来源。\newline\textit{反例}：行业章节全是定性描述；竞争格局只列竞争对手名字 \\
L2-07 & 预测假设有依据 & 评审 & 5 & generator\newline 8 文稿 & 关键预测假设（量、价、毛利率、费用率等对利润影响最大的几项）逐项写出取值和依据（历史趋势、公司指引、行业数据之一）。\newline\textit{反例}：假设表只有数字没有依据；最重要的假设没有解释 \\
L2-08 & 估值方法与参数 & 评审 & 3 & generator\newline 8 文稿 & 说明为什么用这种估值方法，WACC、永续增长率等关键参数的取值依据可在报告中找到。\newline\textit{反例}：只给结果不讲方法；WACC 无拆分或依据 \\
L2-09 & 估值交叉验证 & 评审 & 4 & generator\newline 8 文稿 & 主估值方法之外，至少用一种方法或指标交叉验证（可比公司、隐含倍数、历史估值区间等），并解释差异来源。\newline\textit{反例}：只有 DCF 一个数；列了可比公司但不和 DCF 结果对照 \\
L2-10 & 情景与敏感性 & 评审 & 3 & generator\newline 8 文稿 & 给出乐观、悲观情景各自的关键假设差异与对应价值，并指出价值对哪个变量最敏感。\newline\textit{反例}：只给三个价格不说假设差在哪；敏感性表无解读 \\
L2-11 & 风险落到假设 & 评审 & 4 & generator\newline 8 文稿 & 主要风险能对应到某个预测假设，并给出影响方向或量级（例如对每股价值或利润的影响）。\newline\textit{反例}：风险提示是通用模板（宏观、政策、竞争加剧）；风险与预测假设无关联 \\
L2-12 & 跟踪指标与失效条件 & 评审 & 2 & generator\newline 8 文稿 & 给出未来可验证的跟踪指标或催化剂，或写明结论在什么条件下不成立。\newline\textit{反例}：没有任何前瞻可验证的指标 \\
\midrule
L3-01 & 故事线连贯 & 评审 & 4 & reviewer\newline 9 复核改写 & 各章节围绕核心结论展开，从公司怎么赚钱、发生了什么、往哪里去到值多少钱，有清晰主线；没有与主线无关的章节。\newline\textit{反例}：章节像资料堆砌；结论与正文论证对不上 \\
L3-02 & 事实、推论、观点可区分 & 评审 & 3 & reviewer\newline 9 复核改写 & 读者能分清哪些是披露的事实、哪些是推论、哪些是作者观点（例如用“主要系”“我们预计”“我们认为”），且不靠“事实层面”这类行内标签。\newline\textit{反例}：预测写成既成事实；通篇“事实/推论/观点”标签 \\
L3-03 & 结论先行 & 评审 & 3 & reviewer\newline 9 复核改写 & 多数章节和段落先给判断，再给支撑数据。\newline\textit{反例}：段落先堆数字，最后才给结论或没有结论 \\
L3-04 & 无机器腔 & 评审 & 3 & reviewer\newline 9 复核改写 & 没有明显的模型写作痕迹：自证圆场（“与年报一致”“相互印证”）、程序口吻（“程序预测”“经护栏核对”）、否定对照堆砌（“不是……而是……”）、交代缺失（“年报未披露……只能间接验证”）。只看正文段落：图表来源、表注、免责声明、附录里的方法说明不算。偶发的否定对照或个别程序用语可以通过；某一类痕迹在正文里成了习惯写法才判不通过。\newline\textit{反例}：多个章节的段落都用“不是……而是……”起句；正文段落反复出现“经程序核对”“与年报一致”之类的自证 \\
L3-05 & 数字密度 & 评审 & 2 & reviewer\newline 9 复核改写 & 正文不逐行念表；关键数字突出，读者能看出每段最重要的一两个数。一个数字连同它的同比、占比或上年值算一个数据点；单句串起四五个以上互不相关的数据点才算堆砌。\newline\textit{反例}：一句话依次报出多个科目或多个季度的金额与增速，像在念表；段落数字很多，却没有一个被点明为判断的依据 \\
L3-06 & 图表服务论点 & 评审 & 3 & reviewer\newline 10 排版 & 图表有标题与来源，与相邻正文的论点对应，读者能从图上看出正文说的结论。\newline\textit{反例}：图表与正文无关；图表缺来源 \\
L3-07 & 假设与局限透明 & 评审 & 2 & reviewer\newline 10 排版 & 报告中能找到估值假设汇总与口径简化说明（例如哪些项目未纳入、模型的适用范围）。\newline\textit{反例}：没有任何关于模型局限的说明 \\
\midrule
N-01 & 无来源的外部数据 & 评审 & $-$8\newline 封顶70 & generator\newline 1 资料搜集 & 命中条件：正文出现具体的外部数据或排名（市场规模、市占率、行业排名、价格指数等），既没有 \texttt{[S编号]} 也不属于公司年报披露。命中一处即判命中。\newline\textit{反例}：“公司全球市占率第一”且无来源 \\
N-02 & 数字前后矛盾 & 程序\newline\code{consistency} & $-$8 & generator\newline 8 文稿 & 命中条件：评级、某一情景的每股价值、或某一预测年度的公司整体收入 / 利润 / 每股收益 / 毛利率及其增速，在正文两处给出不同数值（差异超出展示精度）。外部来源、条件句、敏感性与区间表述、分部口径不参与比较。\newline\textit{反例}：摘要评级“买入”、投资建议一节“卖出”；摘要基准每股价值 54.87 元、估值一节 49.88 元 \\
N-03 & 聚合数构成描述错误 & 评审 & $-$6 & generator\newline 8 文稿 & 命中条件：正文对某个合计数（股权价值桥加回项、利润变动合计、收入构成等）的构成描述与模型产物中的拆分不符。请对照 model/ 下的产物核对。\newline\textit{反例}：把加回项说成“财务公司拆出资金”，而产物拆分显示是现金与金融资产 \\
N-04 & 结论前后冲突 & 评审 & $-$8 & generator\newline 8 文稿 & 命中条件：两处结论或判断方向互相冲突且口径相同（例如一处说毛利率改善、另一处说下滑；一处说景气上行、另一处说已见顶）。只看判断与方向，不看数值是否一致（数值由程序另行核对）。\newline\textit{反例}：核心观点说“盈利能力持续改善”，业绩回顾说“毛利率同比下滑、盈利承压”，指同一年同一口径 \\\end{longtable}
}

\paragraph{版本差异。}
v0（2026-10-03）与 v2 的差异：L1-10、L1-13、L3-04、L3-05 四条的判定标准在 v1（2026-10-04）修订，修订前后对照见表~\ref{tab:eval-v1}；
v2（2026-10-05）把 N-02 由评审判定改为程序判定（\code{check} $=$ \code{consistency}），判定标准从“前后矛盾”收窄为“数字前后矛盾”，
并新增评审条目 N-04“结论前后冲突”。v0 与 v1 各 40 条，v2 为 41 条。三个版本的全局参数相同。

\section{劣化算子定义}
\label{app:degrade}

劣化算子定义在 \code{finagent/evaluation/calibrate.py} 的 \code{DEGRADATIONS} 字典中，只改文稿 \file{narrative.md}，全部为确定性文本变换。
文稿以 \code{<!-- section: 名称 -->} 标记分节（排版脚本按节名取用），算子依赖以下辅助函数：

\begin{table}[h]
\centering
\caption{劣化算子使用的辅助函数。}
\label{tab:app-degrade-helpers}
\footnotesize
\begin{tabularx}{\linewidth}{@{}l X@{}}
\toprule
函数 & 行为 \\
\midrule
\code{\_sections(text)} & 返回 $[(\text{节名}, a, b)]$：正文为本节标记之后、下一个标记之前的文字 \\
\code{\_replace\_section(text, name, body)} & 把名为 \code{name} 的所有节的正文替换为 \code{body}（为空即清空，标记保留） \\
\code{\_drop\_headings(text, pattern)} & 删掉标题匹配正则 \code{pattern} 的 \texttt{\#\#\#} 小节，直到下一个 \texttt{\#\#\#} 或节标记 \\
\code{\_hollow\_sections\_titled(text, prefix, pattern)} & 节名以 \code{prefix} 开头且节内第一个 \texttt{\#\#\#} 标题匹配 \code{pattern} 时掏空该节：保留标题、表格、注释与段落数，正文段落与列表项换成“本段从略。” \\
\code{\_paragraphs(text, every, edit)} & 对正文段落（非标题、列表、表格、注释，且不在 title / subtitle / summary 节）每隔 \code{every} 段执行一次 \code{edit} \\
\code{\_flip\_rating(text)} & 只在 summary 节内把第一处“给予‘X’评级”按映射换成相反方向 \\
\code{\_insert\_unsourced(text)} & 在 deep\_industry 节第一段之后插入一段固定的无来源数字段落 \\
\code{degrade(text, kind)} & 执行算子；若文稿未发生任何变化则抛 \code{ValueError}（找不到对应小节） \\
\bottomrule
\end{tabularx}
\end{table}

\begin{table}[h]
\centering
\caption{七个劣化算子的精确定义。}
\label{tab:app-degrade}
\footnotesize
\begin{tabularx}{\linewidth}{@{}l l X@{}}
\toprule
算子 & 目标 & 操作 \\
\midrule
\code{no\_cross\_check} & L2-09 & \code{\_drop\_headings(t, "交叉验证|相对估值")} \\
\code{no\_expectation\_gap} & L2-02 & 依次执行：\code{\_drop\_headings(t, "市场关注|预期差")}；\code{\_hollow\_sections\_titled(\(\cdot\), "deep\_topic", "隐含")}；\code{\_replace\_section(\(\cdot\), "deep\_expectations", "")} \\
\code{generic\_risks} & L2-11 & risks 节替换为四行模板“宏观经济波动风险 / 行业竞争加剧风险 / 政策变化风险 / 技术迭代风险”；并把全文“\textbf{风险提示：}”所在行改为“宏观经济波动；行业竞争加剧；政策变化；技术迭代。” \\
\code{no\_tracking} & L2-12 & \code{\_replace\_section(t, "tracking", "")} \\
\code{unsourced\_number} & N-01 & 在 deep\_industry 节第一段后插入：“\textbf{市场空间测算。}据行业测算，2030 年全球市场规模将达到 4,800 亿元，2025—2030 年复合增速约 62\%，其中中国市场占比约 45\%，公司有望占据其中 18\% 的份额。” \\
\code{contradiction} & N-02 & \code{\_flip\_rating}：映射 卖出$\to$买入、减持$\to$增持、买入$\to$卖出、增持$\to$减持、持有$\to$买入、中性$\to$买入；只改摘要第一处 \\
\code{machine\_tone} & L3-04 & \code{\_paragraphs(t, 2, ...)}：每两段正文的第二段前加“根据程序计算结果，”，后加“（以上数据已通过数字护栏校验，由模型自动生成。）” \\
\bottomrule
\end{tabularx}
\end{table}

\paragraph{变体运行。}
\code{prepare\_variant} 把原工作区整体复制到 \file{\textasciitilde/finagent-evals/calibration/<公司>/<变体>/ws/}（排除 \file{\_\_pycache\_\_} 与绘图缓存），
写入劣化后的文稿，再用该工作区自己的 \file{build\_report.py} 按原运行目录重新渲染 PDF；渲染失败或算子未改动文稿时跳过该变体并打印原因。
原稿变体（\code{clean1}～\code{clean}$k$）不复制、不改动、不重排，直接评审原工作区。每个变体的元数据 \file{variant.json} 记录
\code{kind}、\code{rep}、\code{targets}、\code{desc}。评审结果含 \code{judge\_problems}（未写出判定、异常退出）时整个结果目录被删除，下次运行重跑。

\paragraph{测试。}
\file{tests/test\_calibrate.py} 用一份含全部节标记的最小文稿验证：每个算子都改动文稿；\code{no\_cross\_check} 删去交叉验证而保留 DCF 估值小节；
\code{no\_expectation\_gap} 删去市场关注、反向 DCF 专题与预期差正文而保留核心观点；\code{contradiction} 后全文恰有一处“买入”；
其余算子的目标内容消失或出现。另有测试验证缺少对应节时 \code{degrade} 抛错、汇总按原稿多数判定计算命中与连带改判。

\section{数据结构}
\label{app:schemas}

模型会话写下的抽取与假设文件、外部来源登记，以及评测的输入输出，均为 UTF-8 JSON。下表的字段取自仓库中的真实文件
（\file{data/extractions/贵州茅台\_2025/}、\file{data/extractions/星河智造\_2024/claims.json}）、skill 中的格式说明
（\file{.claude/skills/finagent-research/references/extraction.md}、\file{claims.md}）与源码（\file{finagent/output/sources.py}、\file{finagent/evaluation/}）。

\subsection*{C.1 \file{facts.json}：财务事实（分析师会话写，\code{verify-facts} 校验）}

\begin{table}[h]
\centering\footnotesize
\begin{tabularx}{\linewidth}{@{}l l X@{}}
\toprule
字段 & 必填 & 说明 \\
\midrule
\code{company}, \code{fiscal\_year}, \code{source\_file}, \code{notes} & 顶层 & 公司、年度、年报相对路径、口径说明 \\
\code{facts[]} & 是 & 事实列表，每项字段如下 \\
\quad\code{metric\_id} & 是 & 指标 ID，只能取 \code{finagent schema} 输出的清单（如 \code{revenue}、\code{np\_parent}） \\
\quad\code{period} & 是 & 年度 \code{2024}；半年度 \code{2024H1}；单季 \code{2024Q3}；前三季度累计 \code{2024Q1-3} \\
\quad\code{raw\_value} & 是 & 与原文逐字一致的字符串：保留千分位与小数位，括号负数原样，百分比不带 \% \\
\quad\code{unit} & 是 & 元 / 千元 / 万元 / 百万元 / 亿元；每股收益 元/股；比率 \%；股本 股 / 万股 / 亿股 \\
\quad\code{page} & 是 & PDF 页码（从 1 开始），程序在该页按数字边界回查 \\
\quad\code{label}, \code{table} & 建议 & 原文科目名、所在表名 \\
\quad\code{scope} & 是 & 固定 \code{consolidated}（母公司报表不抄） \\
\quad\code{source\_file} & 多文件时 & 季报等补充文件名，同时用 \code{--supplement} 传入 \\
\bottomrule
\end{tabularx}
\end{table}

\code{verify-facts} 输出 \file{facts.verified.json}，顶层为 \code{facts}、\code{issues}、\code{stats}、\code{rejected}、\code{missing}、\code{check}；
其中 \code{check} 含 \code{ok}、\code{stats}（\code{rule\_facts}、\code{llm\_facts}、\code{merged}、\code{rejected}、\code{extractor}、\code{supplements}）、
\code{rejected}、\code{conflicts}、\code{tieout\_failures}、\code{missing}。阶段 2 闸门与评分条目 L1-01 读的都是 \code{check}。
茅台 2025 的实际值为：规则通道 98 条、模型抽取 98 条、合并 116 条，拒收、冲突、勾稽失败、缺失均为 0。

\subsection*{C.2 \file{extras.json}：经营数据（\code{verify-extra} 逐格校验）}

\begin{table}[h]
\centering\footnotesize
\begin{tabularx}{\linewidth}{@{}l X@{}}
\toprule
字段 & 说明 \\
\midrule
\code{tables[]} & 原表逐字抄录：\code{id}、\code{title}、\code{page}、\code{columns[]}（\code{name}、\code{unit}）、\code{rows[]}（\code{label}、\code{values[]}，与列一一对应，“不适用”或空白写 \code{null}；可加行级 \code{page}、\code{alias}）、可选 \code{note}。金额列单位换算为元，其他单位（\%、pp、吨、元/吨等）原值保存，外币金额拒收 \\
\code{ratios[]} & 派生比率由程序算：\code{id}、\code{label}、\code{num}、\code{den}（引用 \code{[表 id, 行名, 列名]}、\code{["truth", 指标, 期间]} 或 \code{["ratio", 比率 id]}）、\code{unit}、可选 \code{op}（缺省除法；\code{growth} 同比；\code{diff} 差值）、\code{scale}、\code{alias} \\
\bottomrule
\end{tabularx}
\end{table}
茅台 2025 的 \file{extras.json} 含 37 张表、93 个派生比率，例如 \code{price\_moutai}（茅台酒吨价 $=$ 分产品营业收入 $\div$ 销量，单位 万元/吨）。

\subsection*{C.3 \file{sources.json}：外部来源登记（\code{verify-sources} 校验）}

\begin{table}[h]
\centering\footnotesize
\begin{tabularx}{\linewidth}{@{}l X@{}}
\toprule
字段 & 说明 \\
\midrule
\code{questions[]} & \code{id}（Q 编号）、\code{question}、\code{status}（\code{found} / \code{not\_found}）、\code{queries[]}（实际检索词）、\code{sources[]}（来源 id）、\code{phase}（\code{collect} 资料搜集 / \code{followup} 补充调研，缺省视为后者）、\code{category}（资料搜集阶段必填，七类之一：latest、market、consensus、peers、industry、competition、events） \\
\code{sources[]} & \code{id}（须为 S + 数字，唯一）、\code{title}、\code{publisher}、\code{url}（须以 http(s) 开头，只接受实际打开过的网页）、\code{published}（YYYY-MM[-DD]，可选）、\code{accessed}（YYYY-MM-DD）、\code{facts[]} \\
\quad\code{facts[]} & \code{quote}（原文摘句，$\le$ 120 字，必须包含所登记的数字）、\code{value}、\code{unit}、\code{label} \\
\bottomrule
\end{tabularx}
\end{table}
正文引用外部数字时在句末写 \texttt{[S编号]}；数字护栏把 \code{facts[].value} 当作出处，并要求引用它的句子带对应编号。

\subsection*{C.4 \file{claims.json}：研报断言（纠错核查）}

\begin{table}[h]
\centering\footnotesize
\begin{tabularx}{\linewidth}{@{}l X@{}}
\toprule
字段 & 说明 \\
\midrule
\code{report\_file}, \code{extracted\_by} & 顶层：被核查的研报文件、抽取者说明 \\
\code{claims[].metric\_id} & 取自 \code{schema} 的 \code{claim\_metrics}；对不上的断言不写 \\
\code{claims[].period} & 断言所说的期间；未写年份按研报主报告期 \\
\code{claims[].claim\_type} & \code{level} 绝对值、\code{yoy} 同比增速、\code{ratio} 比率、\code{pp\_change} 比率同比变动 \\
\code{claims[].value}, \code{unit} & 原文数字字符串（下降写负数）与原文单位（亿元 / 万元 / 元 / \% / 元/股 / pp） \\
\code{claims[].is\_forecast} & 预测值为 true，程序判“无法核实”，不计为错误 \\
\code{claims[].quote} & 研报原文中连续出现、包含该数字的一段，程序逐字核对 \\
\bottomrule
\end{tabularx}
\end{table}

\subsection*{C.5 \file{valuation\_assumptions.json}：估值假设（观点 / 假设，非财报事实）}

\begin{table}[h]
\centering\footnotesize
\begin{tabularx}{\linewidth}{@{}l X@{}}
\toprule
字段 & 说明 \\
\midrule
\code{status} & 声明全部为观点 / 假设（如“草稿（Claude 起草，待用户确认）”） \\
\code{ticker}, \code{industry}, \code{valuation\_date}, \code{as\_of\_date} & 代码、行业、估值日、报表基准日 \\
\code{current\_price}, \code{price\_source} & 参考股价与来源（外部数据，阶段 6 闸门要求非空，依据须引用 \texttt{[S编号]}） \\
\code{rationale} & 各项假设的文字依据：\code{forecast}、\code{wacc}、\code{terminal\_value}、\code{non\_operating}、\code{dividend} 等 \\
\code{wacc} & \code{risk\_free\_rate}、\code{beta}、\code{equity\_risk\_premium}、\code{size\_premium}、\code{company\_specific\_premium}、\code{country\_risk\_premium}、\code{pre\_tax\_cost\_of\_debt}、\code{marginal\_tax\_rate}、\code{target\_debt\_pct}、\code{target\_equity\_pct}、\code{preferred\_pct}、\code{pre\_tax\_cost\_of\_preferred}（小数） \\
\code{terminal\_value} & \code{method}（如 \code{perpetual\_growth}）、\code{perpetual\_growth\_rate}、\code{exit\_ebitda\_multiple} \\
\code{non\_operating\_assets[]} & 股权价值桥加回 / 扣减项：\code{label}、\code{raw\_value}、\code{unit}、\code{page}、\code{sign}（$+1$ 加回，$-1$ 扣减）；逐项回查年报 \\
\code{scenarios.\{base,downside,upside\}} & \code{description}、\code{revenue\_growth[]}（由分部预测覆盖）、\code{capex\_percent\_revenue[]}、\code{terminal\_growth\_rate}、\code{wacc\_adjustment} \\
\code{three\_statement} & \code{ar\_days\_shift}（分情景）、\code{dividend\_payout} \\
\code{sensitivities} & \code{wacc\_delta}、\code{terminal\_growth\_delta}、\code{exit\_multiple\_delta}、\code{revenue\_growth\_delta}、\code{ebit\_margin\_delta} 五组网格 \\
\code{segment\_forecast} & \code{note}；\code{segments[]}（\code{id}、\code{name}、\code{product\_row}、\code{volume\_row}、分情景的 \code{volume\_growth[]}、\code{price\_growth[]}、\code{gm\_change\_pp[]}）；\code{expense\_ratios}（如 \code{selling} 分情景逐年费用率） \\
\bottomrule
\end{tabularx}
\end{table}

\subsection*{C.6 评测输入输出}

\begin{table}[h]
\centering\footnotesize
\begin{tabularx}{\linewidth}{@{}l X@{}}
\toprule
文件 & 结构 \\
\midrule
\file{rubric.json}（盲评目录） & 列表，每项 \code{id}、\code{title}、\code{criterion}、\code{fail\_examples}（可无）、\code{negative}（布尔） \\
\file{judgments.json} & \code{items[]}：\code{id}、\code{passed}（布尔；负向条目 true 表示命中）、\code{evidence}（$\le$ 80 字原文摘录并标位置，无则“无”）、\code{reason}；\code{summary}（三五句总评） \\
\file{checks.json} & \code{\{条目 id: \{passed, detail, evidence\}\}}，只含程序条目 \\
\file{judge\_run.json} & \code{model}、\code{exit\_code}、\code{started\_at}、\code{elapsed\_sec}、\code{turns}、\code{cost\_usd}、\code{input\_tokens}、\code{output\_tokens}、\code{blind\_dir} \\
\file{score.json} & \code{rubric}、\code{total}、\code{raw}、\code{deduction}、\code{caps[]}（\code{cap}、\code{by}）、\code{layers}（各层 \code{name}、\code{earned}、\code{possible}）、\code{deliverable}、\code{not\_deliverable\_reasons[]}、\code{missing[]}、\code{items[]}（\code{id}、\code{layer}、\code{title}、\code{judge}、\code{points}、\code{passed}、\code{got}、\code{route}、\code{stage}、\code{evidence}、\code{reason}）、\code{judge\_problems[]}、\code{judge\_summary} \\
\file{issues.json} & 列表，每项恰为 \code{id}（I01…）、\code{route}、\code{stage}、\code{severity}（blocker / major / minor）、\code{description}、\code{evidence} \\
\file{variant.json}（校准） & \code{kind}、\code{rep}、\code{targets[]}、\code{desc} \\
人工质检导出 & \code{schema\_version}（\code{finagent.manual\_review.v2}）、\code{company}、\code{rubric\_version}、\code{selection\_policy}、\code{source\_sha256}、\code{baseline\_total}、\code{flagged\_count}、\code{machine\_excel}、\code{review\_sessions[]}（\code{reviewer}、\code{exported\_at}、\code{overall\_comment}、\code{items[]}：\code{item\_id}、\code{layer}、\code{model\_flagged}、\code{human\_reviewed}、\code{human\_verdict}、\code{feedback}） \\
\bottomrule
\end{tabularx}
\end{table}

判定文件的示例（ASCII 化）：
\begin{lstlisting}
{"items": [
  {"id": "L2-01", "passed": true,  "evidence": "... (summary, p.1)", "reason": "..."},
  {"id": "N-01",  "passed": false, "evidence": "none", "reason": "..."}],
 "summary": "..."}
\end{lstlisting}

\section{十一道闸门}
\label{app:gates}

闸门由 \code{report\_pipeline.status(ws)} 逐项判定（正文见第~\ref{sec:pipeline}~节）。阈值常量定义在 \file{finagent/report\_pipeline.py} 顶部（表~\ref{tab:app-gate-const}）；
表~\ref{tab:app-gates} 给出每道闸门的完整判定条件。路径中的 \code{<work>} 为 \file{data/extractions/<公司>\_<年度>}。

\begin{table}[h]
\centering
\caption{闸门阈值常量。}
\label{tab:app-gate-const}
\footnotesize
\begin{tabular}{@{}l r l@{}}
\toprule
常量 & 值 & 含义 \\
\midrule
\code{MIN\_COLLECT\_QUERIES} & 20 & 资料搜集阶段检索次数下限 \\
\code{MIN\_RECENT\_SOURCES} & 3 & 近期来源数下限 \\
\code{RECENT\_DAYS} & 90 & “近期”的天数 \\
\code{MIN\_TOPICS} & 3 & 深度专题数（框架中）与专题笔记份数下限 \\
\code{MIN\_TOPIC\_CHARS} & 1500 & 每份专题笔记去空白后的字数下限 \\
\code{MIN\_TOPIC\_SOURCES} & 3 & 每份专题笔记的独立 \texttt{[S编号]} 数下限 \\
\code{MIN\_FOLLOWUP} & 8 & 补充检索问题数下限 \\
\code{MIN\_NARRATIVE\_CHARS} & 15000 & 文稿正文字数下限 \\
\code{MIN\_DEEP\_SHARE} & 0.6 & 深度研究（\code{deep\_} 开头的节）占比下限 \\
\code{MIN\_PAGES}, \code{MAX\_PAGES} & 25, 32 & PDF 页数区间（按 preview 页图计） \\
\code{MIN\_FIGURES} & 60 & 图表数下限（按 HTML 中“图表 $n$”编号计） \\
\bottomrule
\end{tabular}
\end{table}

{\footnotesize
\begin{longtable}{@{}p{0.13\linewidth}p{0.11\linewidth}p{0.70\linewidth}@{}}
\caption{十一道闸门的判定条件。}\label{tab:app-gates}\\
\toprule
阶段 & 会话 & 通过条件 \\
\midrule
\endfirsthead
\toprule
阶段 & 会话 & 通过条件 \\
\midrule
\endhead
\bottomrule
\endfoot
1 资料搜集 & 研究助理 & \file{<work>/research/sources.json} 与 \file{dossier.md} 都存在；\code{phase=collect} 的问题覆盖全部七个类别；这些问题的检索词合计 $\ge 20$；发布日期在 \code{as\_of} 前 90 天内的来源 $\ge 3$；\code{verify\_sources} 无问题 \\
2 财务抽取 & 分析师 & \file{facts.json} 存在；\file{facts.verified.json} 存在且修改时间不早于 \file{facts.json}；其 \code{check.ok} 为真（拒收、冲突、勾稽失败均为 0） \\
3 分析框架 & 分析师 & \file{report/framework.md} 存在，含“外部调研清单”字样，且“\#\# 深度专题”一节下编号或列表项 $\ge 3$ \\
4 深度调研 & 分析师 & \file{sources.json} 存在；\file{research/topics/*.md} $\ge 3$ 份，每份去空白 $\ge 1500$ 字且独立 \texttt{[S编号]} $\ge 3$；补充检索问题 $\ge 8$；\file{sources.json} 修改时间不早于 \file{framework.md}（确实按框架补查）；\code{verify\_sources} 无问题 \\
5 经营数据 & 分析师 & \file{extras.json} 存在，\code{verify-extra} 拒收 0 格 \\
6 估值假设 & 分析师 & \file{valuation\_assumptions.json} 存在；\code{current\_price} 与 \code{price\_source} 非空；文件中至少一处 \texttt{[S编号]} \\
7 估值运行 & 分析师 & 存在成功的 valuation 运行；其开始时间晚于 facts / extras / 假设 / 可比公司表中最晚的修改时间；带了 \code{--extras}（若有 extras）；若资料搜集写了 \file{research/peers.csv} 且 build 时未给可比公司，运行带了该 \code{--peers} \\
8 文稿 & 主笔 & \file{narrative.md} 存在；运行目录中的 \file{commentary.md} 与当前文稿逐字相同（guard 核对的是当前稿）；\file{commentary\_guard.json} 为 \code{ok}、带了 \code{--sources} 且晚于 \file{sources.json}；正文 $\ge 15{,}000$ 字且深度研究占比 $\ge 60\%$ \\
9 复核改写 & 复核编辑 & \file{narrative.draft.md}、\file{edit\_log.md}、\file{review.md} 都存在；\code{numlock(draft, narrative)} 通过（改写稿的数字与来源编号是初稿的子集） \\
10 排版 & 排版设计师 & 运行目录下 \file{report/report.pdf} 存在且晚于文稿与 \file{build\_report.py}；另存了按 \code{output\_name} 命名的 PDF 且与主文件摘要相同；preview 页数在 25～32；图表 $\ge 60$；\code{check\_render\_record} 对当前 HTML、PDF、排版代码的宽度与字号检查通过 \\
11 交付说明 & 排版设计师 & \file{report/delivery.md} 存在 \\
\end{longtable}
}

\paragraph{阻断规则。}
闸门按顺序判定，任何一道未通过，其后所有“通过”都被改判为未通过并在详情后注明“（前序阶段未通过）”。\code{status} 的 \code{next} 字段给出第一道未通过闸门的下一步动作。
\code{run\_agent} 按会话分组推进：会话负责的最后一道闸门（1、7、8、9、11）已通过则跳过该会话；会话结束后其闸门仍未全部通过则停止整个流水线。

\section{角色启动提示要点}
\label{app:prompts}

每个角色由编排器以一次 \code{claude -p} 无头会话启动，启动提示（prompt）很短，详细的工作说明在工作区 \file{CLAUDE.md}（由 \code{render\_guide} 生成）与 skill 文档中。
生成侧会话命令为 \code{claude -p <PROMPT> --permission-mode acceptEdits --allowedTools <白名单> --output-format stream-json --verbose [--model <模型>]}。

{\footnotesize
\begin{longtable}{@{}p{0.15\linewidth}p{0.12\linewidth}p{0.67\linewidth}@{}}
\caption{各角色的启动提示要点（\file{finagent/report\_pipeline.py}、\file{finagent/evaluation/judge.py}）。}\label{tab:app-prompts}\\
\toprule
提示常量 & 角色 / 模型 & 要点 \\
\midrule
\endfirsthead
\toprule
提示常量 & 角色 / 模型 & 要点 \\
\midrule
\endhead
\bottomrule
\endfoot
\code{RESEARCH\_PROMPT} & 研究助理 / Opus & 只做阶段 1：先读 \file{finagent-report/references/research.md} 第一部分，用 WebSearch / WebFetch 把年报之后的最新信息搜齐，登记 \file{sources.json}、写 \file{dossier.md}；不做财务抽取、不写框架；结束运行 \code{status}，阶段 1 通过即停 \\
\code{ANALYST\_PROMPT} & 分析师 / Opus & 先读 \file{dossier.md}，完成阶段 2～7（抽取、框架含深度专题、深度调研、经营数据、估值假设、估值运行）；不写文稿；每阶段结束跑 \code{status}，阶段 7 通过即停 \\
\code{TASK\_PROMPT} & 主笔 / Opus & 先读 \file{framework.md}、\file{dossier.md}、\file{research/topics/} 与运行产物，按阶段 8 写“年报点评 + 深度研究”，多环定量推导、写长写深；guard 与篇幅闸门通过即停 \\
\code{REVIEW\_PROMPT} & 复核编辑 / Sonnet & “没有参与写作”；先读 \file{finagent-review/SKILL.md}，只做阶段 9；通过即停 \\
\code{LAYOUT\_PROMPT} & 排版设计师 / Opus & 文稿已定稿、不改正文；读 \file{html-report.md} 与参考脚本，从模板写本公司排版脚本；每张图服务一个论点，图表 60 张以上、25～30 页，表格涨跌双色，逐页看 preview；完成阶段 10～11 \\
\code{FIX\_PROMPT} & 修订 / Opus & 返工第 2 轮起：按问题单修改对应阶段产物（补数据、改假设则重跑估值，文稿改完重跑 guard），处理不了的写进 \file{review.md}；阶段 1～8 通过即停 \\
\code{REVISE\_PROMPT} & 复核修订 / Sonnet & 读问题单逐条处理：只改措辞、结构与编排，数字锁照常，不覆盖初稿；需补数字或改结论的写进 \file{review.md}；重跑 guard 与数字锁，用现有脚本重新渲染（不改版面），更新交付说明 \\
\code{judge.PROMPT} & 盲评质检员 / Sonnet & “按本目录 CLAUDE.md 评审这份研报，把结果写进 judgments.json 后停止。”指南要点见第~\ref{sec:eval-judge}~节 \\
\end{longtable}
}

所有生成侧提示都要求“命令直接以 \code{./py} 开头，不加 cd、不串联”，因为无头执行时工具白名单按命令前缀匹配。

\paragraph{工具白名单。}
\begin{sloppypar}
生成侧（\code{ALLOWED\_TOOLS}）：\code{Bash(./py -m finagent *)}、\code{Bash(./py -m finagent.report\_pipeline *)}、\code{Bash(./py data/extractions/*)}、\code{Bash(./py -c *)}、
\code{Bash(cp data/extractions/*)}、\code{Bash(ls *)}、\code{Bash(mkdir -p data/*)}、\code{Read}、\code{Write(./data/**)}、\code{Edit(./data/**)}、\code{WebSearch}、\code{WebFetch}。
工作区 \file{.claude/settings.json} 另设禁止规则：不得读取仓库所在目录及其上级目录；不得编辑或写入 \file{finagent/}、\file{.claude/}、\file{report/resources/}、\file{sample.json}、\file{py}、\file{py.cmd}、\file{CLAUDE.md}。
盲评会话：\code{Read(./**)}、\code{Glob}、\code{Grep}、\code{Write(./judgments.json)}。
\end{sloppypar}

\paragraph{工作区 \file{CLAUDE.md} 的五条底线。}
（1）不做算术：增速、占比、差值、估值都由程序算，正文照产物展示值抄；（2）每个数字来自年报原文（带页码）、运行产物或已登记的外部资料，正文必须过 guard，年报之外的事实先检索、登记、再引用；
（3）估值参数是观点 / 假设，须标明，股价须注明来源，评级取程序结果；（4）校验 3 轮仍不过的问题不硬改数字，写进交付说明；（5）五个会话互不继承上下文，各自做完自己的阶段就停止。
此外 \file{CLAUDE.md} 规定了读取边界（不读引擎源码、不读 \file{runs/} 下 JSON 全文、PDF 按页读且每次不超过 20 页）、
“先搜集、后分析”（年报之外的信息一律联网现查，来源优先级：交易所与公司公告 $>$ 政府与行业协会 $>$ 研究机构与数据商 $>$ 主流财经媒体），
以及“拿不准但影响结论的内容写进复核清单，不写进正文”。

\paragraph{四个 skill 的分工。}
\code{finagent-pipeline}（仓库内的编排入口，不进工作区）：搭工作区、起会话、查闸门、收回、测评，自己不写研报。
\code{finagent-research}（分析师）：铁律、定位关键页、读页抽取、校验与修正（最多 3 轮）、经营数据、研报断言、运行计算、写解读并过护栏。
\code{finagent-report}（研究助理、主笔、排版）：不可协商的约束与交给模型判断的部分、资料搜集与深度调研方法、分析框架推导、文稿模板、reportkit 排版。
\code{finagent-review}（复核编辑）：边界（只改措辞与结构、不碰数字、不改结论与评级、保留节标记与引用）、怎么读（先通读回答三个问题，再逐段改，再对账，最后通读）、
要改掉的写法（否定对照、自证圆场、交代缺失、程序口吻、套话、解释腔、事后圆场、行内标签、数字罗列、重要处定性、无来源判断）与产出（改写稿、\file{edit\_log.md}、\file{review.md}）。

\section{第三方组件}
\label{app:thirdparty}

本附录摘自仓库 \file{THIRD\_PARTY.md}。

\paragraph{OpenAI Public Equity Investing 插件。}
来源为该插件 0.1.31 版，原插件标注为“Proprietary – Internal Use Only”。本仓库\textbf{不包含}原插件代码：早期版本曾通过 \file{third\_party/} 原样引入，现已全部删除，改为自研实现；
仍参照的是计算口径、判定阈值、输出字段名与英文代码值，以保证改造前后结果一致。参照关系如表~\ref{tab:app-third}。

{\footnotesize
\begin{longtable}{@{}p{0.36\linewidth}p{0.36\linewidth}p{0.22\linewidth}@{}}
\caption{参照原插件的模块与自研实现。}\label{tab:app-third}\\
\toprule
参照的原插件内容 & 自研实现 & 一致性校验 \\
\midrule
\endfirsthead
\toprule
参照的原插件内容 & 自研实现 & 一致性校验 \\
\midrule
\endhead
\bottomrule
\endfoot
三表模型构建（三表联动、营运资本、PP\&E、债务与现金归集、勾稽、情景与敏感性、模型状态判定） & \file{tasks/valuation/ts\_engine.py} & 浮点逐位相等 \\
情景敏感性生成（三情景目标价与偏度、倍数 / EPS / KPI / 流动性敏感性、跟踪触发表、就绪度） & \file{tasks/valuation/scenario\_tables.py} & 输出逐字节相等 \\
业绩深度分析（beat / miss、业绩预告中值对比、EPS 质量筛查规则） & \file{tasks/financials/earnings.py} & 逐字段相等 \\
FCFF DCF 模型 & \file{tasks/valuation/dcf.py} & 相对误差 $<10^{-9}$ \\
可比公司估值说明（EV 桥、倍数配对、NM、IQR 离群值、P25/P75 区间） & \file{tasks/valuation/comps.py} & — \\
可比公司英文公式模板 & 未沿用；\code{comps.write\_formula\_workbook} 另行生成中文公式工作簿 & 公式独立重算与缓存值相等 \\
目标反解思路（“股价隐含了什么”、可行性分档） & \file{tasks/valuation/reverse\_dcf.py}：自研单变量反向 DCF（网格扫描 + 二分），分档自定 & 仅借方法；\file{test\_reverse\_dcf.py} 代回正向 DCF \\
研报 QC 的重复数字对账规则与问题分类 & \file{tasks/report\_check/consistency.py}（N-02、guard、纠错核查） & 仅借方法；\file{test\_consistency.py} \\
首次覆盖论点框架、论点跟踪工作流 & \file{finagent-report/}\allowbreak\file{references/analysis.md}、\file{sections.md}：中文自写的写作要求 & 文档，无计算 \\
模型审计中的公式与工作簿控制 & \file{evaluation/workbook\_audit.py} 的 \code{formula\_controls} & 仅借方法；\file{test\_workbook\_audit.py} \\
财务归一化的证据标签与来源协议 & \file{output/guard.py} 的 \code{tagged\_truth\_values} / \code{classify\_text}：标签体系自定（中文五类） & 仅借方法；\file{test\_provenance.py} \\
\end{longtable}
}

\begin{sloppypar}
校验方式：\file{tests/golden/} 为改造前用原插件引擎跑样例得到的产物快照，\file{tests/test\_golden.py} 要求自研实现的产物与之完全相等（无需原插件即可运行）；
2026-09-28 因有意修改口径（EBIT 剔除投资收益与公允价值变动、三表经营费用改为轧差、DCF 由三表驱动、comps 增加利润率）重新生成了快照，各引擎本身未改。
\file{tests/test\_parity\_reference.py} 在本地存在原插件时（\code{FINAGENT\_REFERENCE\_PLUGIN}）用样例、边界值和数百组随机输入逐项对照，找不到原插件时跳过。
\end{sloppypar}

\paragraph{开源依赖。}

\begin{table}[h]
\centering\footnotesize
\begin{tabular}{@{}l l l l@{}}
\toprule
名称 & 版本 & 许可证 & 使用范围 \\
\midrule
pdfplumber & $\ge$0.11 & MIT & PDF 文本与表格抽取 \\
pytest & $\ge$8.0 & MIT & 测试 \\
fpdf2 & 2.8 & LGPL-3.0 & 仅开发时生成样例 PDF，非运行依赖 \\
Textual & $\ge$1.0（开发时 8.2.8） & MIT & 仅终端界面 \file{finagent\_tui/}，可选 \\
Rich & 随 Textual（开发时 15.0.0） & MIT & 终端界面的表格与文本渲染 \\
\bottomrule
\end{tabular}
\end{table}
\file{requirements.txt} 另列可选的 matplotlib（$\ge$3.8，研报画图）与 pillow（$\ge$10，预览缩略图）。\file{finagent/} 其余部分只用 Python 标准库（含 xlsx 写出）。
样例 PDF 内嵌 macOS 系统字体 Arial Unicode 的字形子集，仅用于显示样例数据。

\paragraph{大语言模型与执行器。}
Claude（Anthropic）：生成侧默认 Claude Opus 5.5，复核与盲评默认 Claude Sonnet 5.5，经无头 Claude Code 调用，由 \file{finagent/report\_pipeline.py} 逐角色启动；
用于联网检索与登记来源、读年报抽取数字、拆研报断言、起草估值假设、写文稿、复核改写、写排版脚本与逐条评审。所有计算、校验与计分由本项目代码完成。
Claude Code 为外部命令行工具，用户自行安装与登录，不随源码分发。

\paragraph{自主开发部分。}
\file{finagent/} 下全部代码自主开发：追踪与编排、多 Agent 编排器、文档解析、A 股指标字典与抽取、财务分析与异常信号、利润变动归因、业绩预期对比、DCF、三表联动、
情景敏感性、相对估值与公式工作簿、研报核查、报告与 xlsx 输出、评测与测试。其中参照原插件计算规则的模块如上表所列。

\section{复现说明}
\label{app:repro}

\paragraph{环境。}
Python 3.12（\file{scripts/setup.sh} 接受 $\ge$3.10）；引擎只依赖标准库与 pdfplumber。深度研报流水线另需 Node.js 18+、已登录的 Claude Code
（\code{claude auth login} 或 \code{ANTHROPIC\_API\_KEY}）、Chrome / Chromium / Edge（HTML 打印 PDF，可用 \code{FINAGENT\_CHROME} 指定）与中文字体。
模型由 \code{FINAGENT\_GEN\_MODEL}（默认 \code{claude-opus-5-5}）与 \code{FINAGENT\_QC\_MODEL}（默认 \code{claude-sonnet-5-5}）配置；
工作区默认在 \file{\textasciitilde/finagent-samples/}（\code{FINAGENT\_SAMPLES\_DIR}），盲评目录默认在 \file{\textasciitilde/finagent-evals/}（\code{FINAGENT\_EVAL\_DIR}），二者都必须在仓库目录树之外。

\paragraph{准备。}
\begin{lstlisting}
bash scripts/setup.sh              # venv + requirements + Claude Code check + .env + doctor
bash scripts/setup.sh --install-claude   # install Claude Code via npm if missing
powershell -ExecutionPolicy Bypass -File scripts\setup.ps1   # Windows
python -m finagent.doctor          # --live sends one minimal request per model (billed)
\end{lstlisting}
\code{setup.sh} 依次完成：建 \file{.venv} 并安装 \file{requirements.txt}；检查 \code{claude} 命令；缺 \file{.env} 时从 \file{.env.example} 复制；运行 \code{doctor} 自检。
以下命令假定已激活 \file{.venv}；否则 macOS / Linux 用 \file{.venv/bin/python}，Windows 用 \file{.venv\textbackslash Scripts\textbackslash python.exe}。

\paragraph{一、不调用模型即可复现的部分。}
\begin{lstlisting}
python -m pytest -q tests
python eval/run_eval.py
python -m finagent all --input data/samples/annual_report_2024.md --year 2024 \
  --company <sample-company> --assumptions data/samples/valuation_assumptions.json \
  --report data/samples/research_draft.md --peers data/samples/peers.csv \
  --estimates data/samples/estimates.csv --guidance data/samples/guidance.csv
python -m finagent verify-facts --input "report/resources/<company>/<annual-report>.pdf" \
  --year 2025 --facts data/extractions/<company>_2025/facts.json
python -m finagent valuation --input "<annual-report>.pdf" --year 2025 --company <company> \
  --facts .../facts.json --extras .../extras.json --assumptions .../valuation_assumptions.json
python -m finagent.evaluation score --result eval/results/v2_<company>
python -m finagent.evaluation.calibrate report
\end{lstlisting}
（代码清单中的 \code{<sample-company>} 为虚构样例公司“星河智造”，\code{<company>} 为仓库中已收回抽取文件的公司，如贵州茅台、宁德时代、招商轮船；清单只能用 ASCII，故以占位符表示。）

预期输出：
\begin{itemize}
  \item 测试：164 个通过；本地没有原参照插件时跳过 5 个逐项对照测试。
  \item 样例评测（\file{eval/run\_eval.py}，规则通道）：字段抽取 35/35；研报纠错召回率、准确率 100\%，误报 0。样例数据为虚构且内部自洽，只说明流程正确。
  \item 真实年报：用 \file{data/extractions/} 中保存的抽取文件重跑 \code{verify-facts} 与 \code{valuation}，程序重新做逐字证据校验、规则通道比对与全部计算，结果与提交时一致；
        茅台、宁德、招商三家 301 个核验值中规则通道取到且一致 280 个，冲突 0，勾稽 56/56 通过。
  \item 评测：\code{score} 用已保存的 \file{checks.json} 与 \file{judgments.json} 重算，无需模型，得到表~\ref{tab:eval-results} 的分数；
        \code{calibrate report} 由 \file{eval/calibration/} 重新生成校准报告（表~\ref{tab:eval-stability}、\ref{tab:eval-degrade}、\ref{tab:eval-human}）。
\end{itemize}

\paragraph{二、需要模型会话的部分。}
\begin{lstlisting}
python -m finagent.report_pipeline build --company <company> --year 2025 \
  --price <price> --price-source "<source>"
python -m finagent.report_pipeline refine --workspace "<workspace>"   # generate + QC loop
python -m finagent.report_pipeline run --workspace "<workspace>" [--dry-run]
python -m finagent.report_pipeline status --workspace "<workspace>"
python -m finagent.report_pipeline collect --workspace "<workspace>" --label Opus
python -m finagent.report_pipeline scorecard --workspace A --workspace B
python -m finagent.evaluation run --workspace "<workspace>"
python -m finagent.evaluation.calibrate run --workspace "<workspace>" --repeats 3 --jobs 4
\end{lstlisting}
\code{build} 拒绝未登记在 \file{manifest.json} 中的年报（来源与版本未核对，需加 \code{--allow-unverified}），且 \code{--price} 必须同时给 \code{--price-source}。
会话产出受联网检索结果与模型采样影响，不能逐字复现；可复现的是每个会话的完整事件流（\file{.pipeline/agents/*.jsonl}）、运行记录（\file{.pipeline/agent\_runs.json}）、
全部中间产物及其哈希，以及在这些产物上重跑的全部程序计算与程序检查。\code{collect} 把抽取与文稿收回 \file{data/extractions/<公司>\_<年度>/}，
运行目录收回 \file{runs/}，PDF、\file{sheet.xlsx} 与复核清单收回 \file{report/交付/<MMDD> <公司> <标签>/}。

\paragraph{统一命令行。}
\code{pip install -e ".[tui]"} 后可用 \code{finagent}：\code{research}（引擎任务与校验）、\code{report}（研报流水线）、\code{eval} / \code{calibrate}（质检与校准）、
\code{bench} / \code{test} / \code{doctor}、\code{ls} / \code{tui}（终端界面）。

\paragraph{修改样例时。}
改样例年报须改 \file{data/samples/build\_samples.py} 后重新生成（同时生成样例三季报），并同步更新 \file{eval/gold\_financials.json}、\file{eval/labels.json}、样例 PDF（\file{build\_pdf.py}）
与 \file{data/extractions/} 下的抽取文件；改完代码须跑测试与样例评测，样例评测指标下降视为回归。
